\documentclass[10pt]{article}
\usepackage[T1]{fontenc}
\usepackage[utf8]{inputenc}
\usepackage{lmodern}
\usepackage[margin=0.7in]{geometry}
\usepackage{amsmath,amssymb}
\usepackage{microtype}
\usepackage{graphicx}
\usepackage{float}
\usepackage{booktabs}
\usepackage[font=small,labelfont=bf,skip=3pt]{caption}
\captionsetup[figure]{font=footnotesize}
\captionsetup[table]{font=footnotesize}
\setlength{\parskip}{2pt}
\setlength{\parindent}{0pt}
\setlength{\intextsep}{6pt}
\makeatletter
\renewcommand\maketitle{\begin{center}{\large\bfseries\@title\par}\end{center}}
\makeatother

\title{How far apart are AdamW and Muon at equal loss? Constant-loss paths in weight space}

\begin{document}
\maketitle
\vspace{-0.8em}

%ANSWER%

\textbf{Setup.} The models are the 14M-parameter GPTs on TinyStories of Fig.~1 of the checkpoint report: AdamW and
Muon from one initialization, data orders 0 and 1. The six loss levels are the losses of AdamW (order 0) at steps
397, 1440, 3300, 5000, 7500 and 10000, i.e.\ 4\% to 100\% of its training. At each level, every run contributes the
checkpoint whose training loss is closest to the level. Pairs: AdamW vs Muon on the same data order (blue) and, as
the floor, two data orders of one optimizer (grey).

\textbf{Measures.} \emph{Straight line:} the loss along $(1-t)\theta_A + t\theta_B$; its \emph{rise} is the largest
excess over the linear interpolation of the endpoint losses. \emph{Constant-loss path:} $N=10$ points between
$\theta_A$ and $\theta_B$, started on the line, are moved until each has its target loss (interpolated between the
endpoint losses) while the path is kept short. Each pair gets the same budget: a stochastic string method (up to 6000
steps; each step does Adam on $(L_i - c_i)^2$ and pulls the points together), Levenberg--Marquardt steps, curve
shortening within the level set, and, if the points are still off the level, one 2500-step continuation of the
string method without the pull. Every step uses fresh training sequences. A path is \emph{valid at $\tau$} if the loss
on 512 held-out sequences stays within $\tau$ of the target at 9 points per segment; up to four times, segments that
leave this band get a midpoint and the optimization continues. We report $\tau = 0.02$ nats and, for the pairs that
miss it, $\tau = 0.1$: the optimization is the same, only the check and the midpoints use 0.1. \emph{Detour:} path
length$/\|\theta_A - \theta_B\| - 1$. \emph{Prediction-space ratio:} the sum of Jensen--Shannon distances between
the next-token predictions of consecutive models along the path, divided by that between the endpoints.

%RESULTS%

\end{document}
